\scanpage{64}
\deleted{\unsure{de}} $\Omega(D'')$ \deleted{\unsure{et}} et $\Omega(D)$, et $\Omega(D'')$ et $\Omega(D')$, donc
$\Omega(D)$ et $\Omega(D')$, et on gagne. Si par contre
\unsure{toutes|t\textsuperscript{tes}} les psdroites qui ne passent pas par $t$ passent
par $u$, ou si elles passent toutes par $u'$, on est
dans le cas des deux faisceaux de ps. droites,
de sommet $t$ et $u$ (resp. $u'$), qui demande
encore un traitement à part. Mais supposons
qu'il n'existe pas de $D''$ qui ne passe ni par
$u$ ni par $u'$ (ni par $t$, bien-entendu) mais
qu'il en existe \inserted{$\unsure{\Gamma|\Gamma_1}$} qui ne passe pas par $u'$, \inserted{(mais par $u$)} et d'autres qui ne passe par $u^{\deleted{\unsure{\prime}}}$ (mais par $u'$), alors
\figure{\begin{tikzpicture}[x=.013cm,y=-.013cm,every node/.style={font=\scriptsize,inner sep=1pt}]
\draw (68,339) -- (208,272) -- (467,151) -- (596,105);
\draw (59,309) -- (208,272) -- (472,214) -- (606,180);
\draw (70,192) .. controls (191,204) and (349,208) .. (472,214) -- (691,241);
\draw (270,401) -- (420,216) -- (467,151) -- (520,69);
\draw (462,84) -- (467,151) -- (472,214) .. controls (476,327) and (473,411) .. (491,478) -- (498,532);
\draw (121,185) -- (208,272) -- (307,347) -- (491,478) -- (524,494);
\node at (554,44) {$\deleted{\unsure{D}}\Gamma$}; \node at (740,239) {$\Gamma'$};
\node at (624,100) {$D$}; \node at (635,173) {$D'$}; \node at (493,325) {$\Delta$};
\node at (330,395) {$\Delta'$}; \node at (171,256) {$t$};
\node at (452,143) {$u$}; \node at (483,132) {$\unsure{\wedge|x}$}; \node at (490,229) {$u'$};
\node at (512,465) {$s$};
\end{tikzpicture}}
\unsure{on a}
\[
\begin{aligned}
\Omega(\Gamma)&\simeq\Omega(D')\\
\Omega(\Gamma')&\simeq\Omega(D)
\end{aligned}
\]
et on peut décrire de \unsure{m\textsuperscript{e}|$=$} $\Omega(\Gamma)\simeq\Omega(\Gamma')$
(donc $\Omega(D)\simeq\Omega(D')$) \uline{pourvu que} l'intersection \inserted{$r$} de $\Gamma$ et $\Gamma'$
par malheur ne se trouve encore just\unsure{ment|mt}
sur $\Delta'$ (de sorte que $\Gamma$ et $\Gamma'$ font par r. à $s$
et $\Delta,\Delta'$ la m\textsuperscript{e} figure que tantôt $D,D'$).
Dans ce cas, les 6 psdroites $D,D',\Delta,\Delta',\Gamma,\Gamma'$ sont le
% archive sheet number: 63
